% !TeX root = ../navier-stokes-manuscript.tex
\subsection{What a finite last time would require}\label{sub:WhatFiniteFailureRequires}

Assume that the maximal time is finite, and write
\[
  T_*<\infty,
  \qquad
  I_*:=(0,T_*).
\]
The solution is smooth at every time in \(I_*\).
Any failure is pushed entirely into the terminal approach, where a norm capable of restarting the same solution has to lose its bound.

The equation still preserves one decisive low-order quantity all the way to that approach.
Pairing \eqref{eq:NavierStokesSystem} with \(u\), incompressibility removes transport and pressure, while integration by parts turns viscosity into dissipation:
\begin{equation}\label{eq:BodyEnergyIdentity}
  \frac12\norm{u(t)}_2^2
  +\nu\int_0^t\norm{\nabla u(s)}_2^2\,ds
  =\frac12\norm{u_0}_2^2.
\end{equation}
Thus kinetic energy stays finite at every time, and one spatial derivative is square-integrable after time has been accumulated.
Neither statement puts a ceiling over the gradient at each instant.

That missing pointwise-in-time control is where the three-dimensional mechanism lives.
With
\[
  \omega:=\nabla\times u,
  \qquad
  S:=\frac12\bigl(\nabla u+(\nabla u)^{\mathsf T}\bigr),
\]
the vorticity obeys
\begin{equation}\label{eq:VorticityEquation}
  D_t\omega
  :=\partial_t\omega+(u\cdot\nabla)\omega
  =S\omega+\nu\Delta\omega.
\end{equation}
The strain \(S\) can lengthen vorticity in one direction while incompressibility contracts the fluid transversely.
Viscosity acts on the spatial curvature created during that same deformation.
Both effects are generated by the same velocity field at the same time.

A fixed high Sobolev roof would stop the process.
For every integer \(m\geq3\), the standard differentiated estimate is
\begin{equation}\label{eq:HigherSobolevContinuationEstimate}
  \frac12\frac d{dt}\norm{u}_{H^m}^2
  +\nu\norm{\nabla u}_{H^m}^2
  \leq
  C_m\norm{\nabla u}_\infty\norm{u}_{H^m}^2.
\end{equation}
At this level Sobolev embedding controls \(\norm{\nabla u}_\infty\), and a bounded late-time slice generates a new local interval of uniform positive length.
Gronwall's inequality propagates every fixed \(H^m\) norm whenever \(\norm{\nabla u}_\infty\) is integrable, and Proposition~\ref{prop:LocalMaximalHistory} then restarts the solution from the resulting \(H^3\) terminal datum.

The finite-time threat is now concrete.
Kinetic energy can stay bounded while the variation of the velocity is driven across shorter distances and escapes every restart-grade roof.
To see that threat without asking wobble, fragmentation, or lost alignment to rescue the flow, keep one vortex perfectly centred and let the surrounding fluid continue to stretch it.
